\section{Introduction}
\label{sec:introduction}

A central objective of embodied intelligence is to enable robots to perform manipulation tasks in the physical world. Recent advances in generalist robot policies and embodied foundation models have led to increasingly capable systems across diverse manipulation scenarios~\citep{intelligence2026pi,zhang2026joyaira01foundationmodel,li2026causalworldmodelingrobot,li2025simplevla,zhang2026hy,lyu2026lda,yuan2026qwen}. As these systems continue to advance, reliable evaluation becomes essential for understanding their capabilities, limitations, and remaining failure modes. Existing benchmarks have investigated important aspects of robot manipulation, including language-conditioned long-horizon execution~\citep{mees2022calvin}, generalization~\citep{chen2025robotwin}, memory-dependent manipulation~\citep{chen2026rmbench}, precise manipulation~\citep{chen2026univtac,lan2025autobio}, and real-world deployment~\citep{yakefu2025robochallenge,chen2026manipulationnet,atreya2025roboarena}. Despite this progress, current evaluation protocols remain insufficient for systematically assessing generalist manipulation policies. Many benchmarks rely on relatively simple or short-horizon tasks, and task variations are often introduced primarily through changes in objects, layouts, or language expressions. While such variations are useful for measuring distributional robustness, they often preserve similar underlying manipulation patterns, limiting diagnostic coverage over distinct challenges such as generalization, memory, precision, open-semantic grounding, sequential execution, bimanual coordination, and tool use.

Another key limitation is the separation between simulation and real-world evaluation. Simulation-based benchmarks are efficient and scalable, making them suitable for rapid feedback and model iteration; however, they cannot fully capture contact-rich dynamics, actuation errors, perception noise, and other physical factors that emerge during deployment. Real-world evaluation provides a more direct assessment of policy behavior under physical-world conditions, but is costly, time-consuming, and difficult to reproduce without standardized hardware setups, scene reset procedures, evaluation protocols, and deployment interfaces. These limitations motivate a unified benchmark that combines scalable simulation-based diagnosis with reproducible real-world validation.

We introduce \textbf{\textcolor{Blue_1}{RoboDojo}}, a unified sim-and-real benchmark for efficient, comprehensive, and reproducible evaluation of generalist robot manipulation policies. RoboDojo contains 42 simulation tasks and 18 real-world tasks. The simulation benchmark provides broad capability coverage across Generalization, Memory, Long-Horizon, Precision, and Open, and supports heterogeneous parallel evaluation in Isaac Sim, enabling different tasks and scenes to run concurrently. The real-world benchmark complements simulation by exposing policies to challenging physical deployment conditions, including contact-rich interactions, perception noise, actuation errors, and environmental variations. To support reproducible physical testing, we design \textbf{\textcolor{Blue_1}{RoboDojo-RealEval}}, a remote real-world evaluation system that standardizes the hardware setup, workspace layout, lighting condition, scene reset procedure, evaluation protocol, and deployment interface. We further provide \textbf{\textcolor{Blue_1}{XPolicyLab}}, a unified infrastructure for policy development and deployment, enabling policies to be integrated once and evaluated across RoboDojo simulation and real-world settings with minimal policy-side adaptation.

RoboDojo is designed not only to expand benchmark scale, but also to establish a unified evaluation loop for diagnosing and improving generalist robot policies. In simulation, it supports fast policy iteration through diverse task designs and capability-oriented analysis. In the real world, it examines whether policy improvements transfer to challenging physical deployment conditions under standardized and reproducible settings. Together with continuously updated test cases, diagnostic analysis, and a public leaderboard, RoboDojo provides a systematic platform for measuring progress toward robust generalist manipulation. Our contributions are summarized as follows:

\begin{itemize}
\item We develop a unified sim-and-real evaluation system for robot manipulation policies, using a shared policy interface and evaluation pipeline across simulation and real-world testing. The system supports heterogeneous parallel simulation for efficient feedback and RoboDojo-RealEval for standardized, reproducible, and remote physical evaluation.

\item We construct the \textbf{RoboDojo Benchmark}, including 42 simulation tasks and 18 real-world tasks. The simulation tasks cover five capability dimensions: Generalization, Memory, Long-Horizon, Precision, and Open, while the real-world tasks span three robot embodiments and evaluate policies under challenging physical deployment conditions.

\item We build \textbf{XPolicyLab}, a standardized infrastructure for policy development and deployment. XPolicyLab integrates 30 robot manipulation models into a shared framework, enabling policies to be developed, deployed, and evaluated across RoboDojo benchmarks with minimal policy-side adaptation.

\item We conduct extensive evaluations of existing robot manipulation policies on RoboDojo. Based on these results, we establish a public leaderboard and provide systematic analysis of current policy limitations across simulation and real-world settings, highlighting key directions for future generalist manipulation policies.

\end{itemize}

% \section{Introduction}
% \label{sec:introduction}

% A central objective of embodied intelligence is to enable robots to complete manipulation tasks in the physical world. Recent progress in generalist robot policies, Embodied Foundation Models have led to increasingly capable systems across diverse manipulation scenarios~\citep{intelligence2026pi,zhang2026joyaira01foundationmodel,li2026causalworldmodelingrobot,li2025simplevla}. As policy capabilities improve, reliable evaluation becomes essential for understanding their strengths and limitations. Existing benchmarks have studied important aspects of robot manipulation, including language-conditioned long-horizon execution~\citep{mees2022calvin}, generalization~\citep{chen2025robotwin}, memory-dependent manipulation~\citep{chen2026rmbench}, and real-world deployment evaluation~\citep{yakefu2025robochallenge,chen2026manipulationnet,atreya2025roboarena}. Despite these efforts, current evaluation protocols remain insufficient for systematically assessing generalist robot manipulation policies. Many benchmarks rely on relatively simple or short-horizon tasks, and their variations are often constructed mainly by changing objects, layouts, or language expressions. Such variations are useful, but they may still share similar underlying manipulation skills and cover only a limited set of capability dimensions. This makes it difficult to diagnose policy weaknesses across distinct challenges such as generalization, memory, precision, open-semantic grounding, sequential execution, bimanual coordination, and tool use.

% Another key limitation is the separation between simulation and real-world evaluation. Simulation-based benchmarks are efficient and scalable, making them suitable for rapid feedback and model iteration, but they cannot fully capture contact-rich dynamics, actuation errors, perception noise, and other physical factors in deployment. Real-world evaluation provides a more direct test of policy behavior, but it is costly, time-consuming, and difficult to reproduce, especially when the hardware setup, scene reset procedure, evaluation protocol, and deployment interface are not standardized. This motivates a unified benchmark system that connects efficient simulation feedback with reproducible real-world validation, rather than treating simulation and the real world as two independent evaluation settings.

% We introduce \textbf{\textcolor{Blue_1}{RoboDojo}}, a unified sim-and-real benchmark for efficient, comprehensive, and reproducible evaluation of generalist robot manipulation policies. RoboDojo contains 42 simulation tasks and 18 real-world tasks. The simulation benchmark is designed for rapid feedback and model iteration, while providing broad capability coverage across Generalization, Memory, Long-Horizon, Precision, and Open. The real-world benchmark complements simulation by exposing policies to challenging physical-world manipulation conditions, including contact-rich interactions, perception noise, actuation errors, and environmental variations. To accelerate simulation evaluation, RoboDojo implements heterogeneous parallel evaluation in Isaac Sim, enabling different tasks and scenes to run concurrently. For physical testing, we design \textcolor{Blue_1}{\textbf{RoboDojo-RealEval}}, a reproducible real-world evaluation system with remote cloud access. RoboDojo-RealEval standardizes the hardware setup, workspace layout, lighting condition, scene reset procedure, evaluation protocol, and deployment interface, enabling consistent physical-world validation across time and users. We further provide \textcolor{Blue_1}{\textbf{XPolicyLab}}, a unified infrastructure for policy development and deployment, allowing policies to be integrated once and evaluated across RoboDojo simulation and real-world settings with minimal policy-side adaptation.

% RoboDojo is designed not only to increase the number of benchmark tasks, but also to establish a unified evaluation loop for diagnosing and improving generalist robot policies. In simulation, RoboDojo provides diverse task designs and capability dimensions for fast policy iteration. In the real world, RoboDojo evaluates whether policy improvements can transfer to challenging physical deployment conditions under standardized and reproducible settings. Together with continuously updated test cases, diagnostic analysis, and a public leaderboard, RoboDojo provides a systematic platform for measuring progress toward robust generalist manipulation. Our contributions are summarized as follows:

% \begin{itemize}
%     \item We develop a unified sim-and-real evaluation system for robot manipulation policies. The system uses a shared policy interface and evaluation pipeline across simulation and real-world testing. It supports heterogeneous parallel simulation for efficient feedback and RoboDojo-RealEval for standardized, reproducible, and remote physical evaluation.

%     \item We construct the \textbf{RoboDojo Benchmark}, including 42 simulation tasks and 18 real-world tasks. The simulation tasks cover five capability dimensions: Generalization, Memory, Long-Horizon, Precision, and Open. The real-world tasks span three robot embodiments and expose policies to challenging physical deployment conditions.

%     \item We build \textbf{XPolicyLab}, a standardized infrastructure for policy development and deployment. XPolicyLab integrates \XpolicylabNum robot manipulation models into a shared framework, allowing policies to be developed, deployed, and evaluated across RoboDojo simulation and real-world benchmarks with minimal policy-side adaptation.

%     \item We conduct extensive evaluations of existing robot manipulation policies on RoboDojo. Based on these results, we establish a public leaderboard and provide systematic analysis of current policy limitations across simulation and real-world settings, highlighting key directions for future generalist manipulation policies.
% \end{itemize}